\section{Deney Düzeni}
\label{sec:setup}

Resmi test kümesi ancak Ocak 2027'de yayımlanacağı için açık dev kümesini test kümesi olarak kullanıyor ve model seçiminde hiç kullanmıyoruz. Bunun yerine \emph{eğitim} kaynak çiftlerinin \%10'unu doğrulama dilimi olarak ayırıyoruz (her track'te aynı çiftler; böylece diller arası sızıntı olmaz). Bu dilimde öğrenme oranını $\{2,5\}\cdot10^{-5}$ arasından (EN üzerinde) ve en iyi epoch'u (tek dilli eğitimde 10, karma eğitimde 4 epoch içinden) seçiyoruz. Seçim her zaman ayrı çözümlemeyle yapılır, dolayısıyla çözümleyiciden bağımsızdır. AdamW, 32'lik yığın ve \%10 ısınma kullanıyor, dizüstü bilgisayarın 18\,GB belleğine sığmak için kelime gömmelerini donduruyoruz (Ek~\ref{app:hparams}). Her yapılandırma üç tohumla çalıştırılır (Apple M3 Pro'da tek dilli koşu başına 3.5--6 dakika). Tüm skorlar resmi puanlayıcıdan alınmıştır.

\paragraph{Kaynaklar.} Yalnızca açık kodlayıcı ağırlıklarını ve resmi eğitim verisini kullanıyoruz.\footnote{\url{https://inigolopezgazpio.net/SemEval-2027-Task-2-DiCo-NLI/}} Tek harici veri, mDeBERTa-NLI'ın eğitildiği NLI derlemleridir (XNLI, çevrilmiş İngilizce NLI). PhrasIS kullanılmamıştır.

\section{Sonuçlar}
\label{sec:results}

\begin{table*}[t]
\centering\small
\setlength{\tabcolsep}{4pt}
\resizebox{\textwidth}{!}{%
\begin{tabular}{lccccccccccccccc}
\toprule
 & \multicolumn{3}{c}{EN} & \multicolumn{3}{c}{ES} & \multicolumn{3}{c}{EU} & \multicolumn{3}{c}{Karma} & \multicolumn{3}{c}{Ortalama} \\
\cmidrule(lr){2-4} \cmidrule(lr){5-7} \cmidrule(lr){8-10} \cmidrule(lr){11-13} \cmidrule(lr){14-16}
Sistem & F1 & Soft & Hard & F1 & Soft & Hard & F1 & Soft & Hard & F1 & Soft & Hard & F1 & Soft & Hard \\
\midrule
Rastgele & 25.3 & 18.1 & 5.5 & 25.3 & 18.1 & 5.5 & 25.3 & 18.1 & 5.5 & 25.1 & 18.4 & 6.2 & 25.2 & 18.1 & 5.7 \\
Sabit \EQ & 11.0 & 100.0 & 31.4 & 11.0 & 100.0 & 31.4 & 11.0 & 100.0 & 31.4 & 11.0 & 100.0 & 31.4 & 11.0 & 100.0 & 31.4 \\
\midrule
XLM-R & 69.1 & 82.7 & 72.1 & 54.2 & 66.1 & 51.9 & 61.1 & 80.5 & 63.7 & 63.4 & 79.3 & 64.9 & 62.0 & 77.1 & 63.1 \\
\quad + ortak çözümleme & 69.6 & 93.9 & 77.3 & 58.6 & \textbf{93.5} & 68.4 & 61.0 & \textbf{94.3} & 70.0 & 63.8 & 88.4 & 69.7 & 63.3 & \textbf{92.5} & 71.3 \\
\midrule
mDeBERTa & 79.6 & 87.1 & 81.0 & 77.3 & 85.6 & 78.3 & 76.5 & 82.7 & 75.7 & 73.4 & 81.8 & 73.8 & 76.7 & 84.3 & 77.2 \\
\quad + ortak çözümleme & 80.3 & 91.8 & 84.5 & 77.9 & 89.3 & 81.1 & \textbf{76.7} & 87.5 & 78.8 & 73.5 & 86.2 & 76.6 & 77.1 & 88.7 & 80.2 \\
\midrule
mDeBERTa-NLI & 82.8 & 88.9 & 84.6 & \textbf{79.0} & 85.0 & 79.5 & 75.1 & 81.8 & 75.1 & 74.7 & 82.4 & 74.7 & 77.9 & 84.5 & 78.5 \\
\quad + ortak çözümleme & 83.0 & 94.2 & 87.5 & 78.7 & 90.4 & 82.2 & 75.6 & 86.8 & 78.5 & 74.8 & 88.4 & 78.3 & 78.0 & 89.9 & 81.6 \\
\midrule
mDeBERTa + $\mathcal{L}_\text{cons}$ & 80.7 & 90.0 & 83.9 & 76.5 & 86.9 & 78.6 & 75.1 & 85.4 & 77.4 & 72.4 & 85.7 & 75.4 & 76.2 & 87.0 & 78.8 \\
\quad + ortak çözümleme & 80.8 & 92.5 & 85.3 & 76.8 & 91.1 & 81.2 & 75.3 & 88.0 & \textbf{79.3} & 72.5 & 88.6 & 77.2 & 76.3 & 90.0 & 80.8 \\
\midrule
mDeBERTa-NLI, 3 tohum topluluğu & 83.7 & 89.9 & 86.6 & 78.8 & 87.4 & 80.5 & 76.4 & 82.7 & 77.3 & \textbf{76.0} & 83.9 & 76.6 & 78.7 & 86.0 & 80.2 \\
\quad + ortak çözümleme & \textbf{84.3} & \textbf{94.9} & \textbf{89.9} & 78.9 & 91.3 & \textbf{82.7} & 76.2 & 85.9 & 79.1 & 75.9 & \textbf{88.9} & \textbf{79.5} & \textbf{78.9} & 90.3 & \textbf{82.8} \\
\midrule
\textit{mDeBERTa + struct (O2)} & \textit{91.2} & \textit{100.0} & \textit{89.5} & \textit{90.2} & \textit{100.0} & \textit{88.2} & \textit{88.7} & \textit{100.0} & \textit{86.5} & \textit{87.9} & \textit{100.0} & \textit{85.6} & \textit{89.5} & \textit{100.0} & \textit{87.5} \\
\bottomrule
\end{tabular}}
\caption{Dev sonuçları (\%). Öğrenen sistemler: 3 tohumun ortalaması (standart sapmalar Ek~\ref{app:std}'de); \emph{3 tohum topluluğu} satırları ortalama değil, tek bir topluluğun (olasılık ortalaması) sonucudur. \emph{ortak}: ortak çift yönlü çözümleme (Denk.~\ref{eq:joint}). Kalın: en iyi öğrenen (tanı dışı) değer. Son satır O2'yi kullanır; yalnızca tanı amaçlıdır.}
\label{tab:main}
\end{table*}


\paragraph{Kodlayıcılar.} mDeBERTa, XLM-R'ı ortalamada 14.7 F1 puanı geride bırakır (Tablo~\ref{tab:main}). XLM-R bu küçük kümelerde kararsızdır da: bir İspanyolca tohumu 0.41 doğrulama F1'ine çöktü. İngilizcede mDeBERTa (79.6 F1 / 81.0 HardCons), iki düzen birebir aynı olmasa da DeBERTa-v3 pilot temel sisteminin \citep[75 / 79;][]{apaolaza2026assessing} üzerindedir. Skorlar EN'den ES'e, EU'ya ve Karma'ya doğru düşer. Karma track'te en zor kombinasyon İspanyolca--Baskça'dır (en-es için 75.1'e karşı 71.6 F1).

\paragraph{Ortak çözümleme.} 72 değerlendirmenin tamamında (6 sistem $\times$ 3 tohum $\times$ 4 track) ortak çözümleme HardCons'u \emph{her seferinde} artırdı; ortalama artış $+3.9$ puandır. Kazanç mDeBERTa'da $+3.0$, XLM-R'da $+8.2$'dir: model ne kadar zayıf ve kararsızsa kazanç o kadar büyüktür. Ortak çözümleme SoftCons'u da her değerlendirmede ($+6.5$), F1'i ise hafifçe artırdı ($+0.5$; 72 değerlendirmenin 51'inde pozitif). Ortak çözümleme tohum topluluğuyla birbirini tamamlar. En iyi sistemimiz, NLI ön eğitimli mDeBERTa'nın ortak çözümlemeli 3 tohumlu topluluğudur ve ortalama 78.9 F1 ile 82.8 HardCons'a ulaşır. $y^\star=\NO$ olan çiftler tutarsız sayıldığı için SoftCons \%100'ün altında kalır.

\paragraph{NLI ön eğitimi.} NLI ile ince ayar yapılmış modelden başlamak İngilizcede ($+3.2$ F1) ve İspanyolcada ($+1.7$), Karma'da ise hafifçe ($+1.3$) fayda sağlar. XNLI'da bulunmayan bir dil olan Baskçada fayda sağlamaz ($-1.4$). Ortak çözümleme altında ortalama $+1.4$ HardCons ekler.

\paragraph{Eğitimde mi, çözümlemede mi tutarlılık?} Tutarlılık kaybı amaçlandığı gibi çalışır. Tutarsız çiftlerin payını düşürür (\%2.5$\rightarrow$\%1.7) ve ayrı çözümlemede HardCons'u $+1.6$, SoftCons'u $+2.7$ artırır; bedeli $-0.5$ F1'dir. Ancak ortak çözümleme tek başına bu HardCons kazancının yaklaşık iki katını verir ($+3.0$), ortak çözümlemenin altına kaybı eklemek ise yalnızca $+0.6$ getirir (80.2$\rightarrow$80.8). Dolayısıyla tutarlılığı çözümleme anında zorlamak, onu öğretmekten hem daha ucuz hem daha etkilidir.

\begin{table}[t]
\centering\small
\setlength{\tabcolsep}{2.5pt}
\begin{tabular}{@{}llccccc@{}}
\toprule
Model & Çöz. & Doğru & \NO & Tuts. & Yanlış & F1\textsubscript{\NO} \\
\midrule
XLM-R & ayrı & 63.1 & 9.9 & 13.0 & 14.0 & 33.1 \\
 & ortak & 71.3 & 7.5 & 0.0 & 21.2 & 33.0 \\
mDeBERTa & ayrı & 77.2 & 13.2 & 2.5 & 7.1 & 53.9 \\
 & ortak & 80.2 & 11.3 & 0.0 & 8.5 & 54.0 \\
\bottomrule
\end{tabular}
\caption{Altın tersinir çiftlerin sonucu (çiftlerin \%'si; dev, track ve tohum ortalaması): iki yön de doğru (Doğru), bir yöne \NO tahmin edilmiş, tutarsız (Tuts.) ya da tutarlı ama yanlış. F1\textsubscript{\NO}: \NO sınıfının örnek düzeyindeki F1'i.}
\label{tab:errors}
\end{table}


\paragraph{Hatalar nereden geliyor?} Tablo~\ref{tab:errors}, her altın tersinir çiftin sonucunu ayrıştırır. mDeBERTa'da çiftlerin yalnızca \%2.5'i yön açısından tutarsızdır; baskın hata, tersinir bir çift için \NO tahmin etmektir (\%13.2). \NO sınıfı yalnızca 54 F1'e ulaşır. Ortak çözümleme tüm tutarsızlıkları giderir ve bazı hatalı \NO tahminlerini de düzeltir (\%13.2$\rightarrow$\%11.3), çünkü emin olan bir yön diğerini geçersiz kılar. Tablo~\ref{tab:main}'deki tanı amaçlı \emph{struct} satırı, \NO kararının ne kadar önemli olduğunu ölçer: \NO'yu model yerine çift yapısıyla (O2) belirlemek ortalama F1'i 76.7'den 89.5'e çıkarır. Dolayısıyla açık veride görevin en zor kısmı, ilişkinin \emph{hangi yönde} geçerli olduğu değil, \emph{tersinir olup olmadığına} karar vermektir. Test verisi bu yapıyı gizlemedikçe, bu kısım aynı zamanda en kolay kestirme yoldur.

\begin{table}[t]
\centering\small
\setlength{\tabcolsep}{4pt}
\begin{tabular}{@{}lcccc@{}}
\toprule
Eğitim verisi & EN & ES & EU & Karma \\
\midrule
Kendi verisi & 84.5 & 81.1 & 78.8 & 76.6 \\
Yalnız EN & 84.0 & 80.7 & 68.5 & 54.4 \\
Birleşik & 84.8 & 83.6 & 81.5 & 77.7 \\
\bottomrule
\end{tabular}
\caption{Ortak çözümlemeli mDeBERTa'nın eğitim verisine göre dev HardCons değeri (\%, 3 tohum): track'in kendi verisi, yalnızca Track~1 (diğerlerinde sıfır-örnek) ya da dört track'in birleşimi.}
\label{tab:transfer}
\end{table}


\paragraph{Diller arası aktarım.} Yalnızca İngilizceyle eğitilen bir model İspanyolcada HardCons'unu korur (kendi verisiyle 81.1'e karşı 80.7; Tablo~\ref{tab:transfer}), ancak 3.5 F1 puanı kaybeder. \NO'yu daha seyrek tahmin eder (duyarlılık \%59 yerine \%43); bu, \NO örneklerinde F1'e mal olur ama tersinir çiftleri korur. Baskçada 10, karma track'te 22 HardCons puanı kaybeder; yani diller arası çiftler diller arası denetim gerektirir. Dört eğitim kümesini birleştirmek en iyi tek model ayarıdır (kendi verisiyle 80.2'ye karşı ortalama 81.9 HardCons); en büyük kazanç Baskça ($+2.7$) ve İspanyolcada ($+2.5$).

\section{Çalışma Planı ve Görev Dağılımı}
\label{sec:plan}

Bölüm~\ref{sec:method}'teki tüm bileşenler uygulandı. Final raporu (28 Aralık) için planımız:
(i) NLI ön eğitimi ile birleşik eğitimi 3 tohumlu toplulukta birleştirmek ve \emph{her} örneğin yer değiştirmiş girdisini modelin kendisiyle puanlamak (ters yön dosyada olmasa da);
(ii) \NO logitini doğrulama diliminde kalibre etmek ve ters çevrilmiş \NO eğitim örneklerini (bunlar da \NO'dur) 0.5 kayıp ağırlığıyla eklemek;
(iii) $\lambda\in\{0.1, 0.3, 1\}$ ayarı;
(iv) nitel hata analizi; isteğe bağlı olarak küçük, açık bir LLM ile sıfır-örnek deney (ayrı ve ortak çözümleme).

\begin{table}[h]
\centering\small
\setlength{\tabcolsep}{3pt}
\begin{tabular}{@{}lp{0.74\columnwidth}@{}}
\toprule
Üye & Kod / analiz; plan maddesi; yazdığı bölümler \\
\midrule
Küçük & Puanlayıcı entegrasyonu, temel sistemler, hata analizi (Tablo~\ref{tab:errors}); (iv); \S1, \S6 \\
Şahin & Eğitim altyapısı, ortak çözümleme, deneyler; (i); \S4, 5, 7, 8 \\
Okudurlar & Kodlayıcılar, NLI aktarımı, tutarlılık kaybı; (iii); \S3, \S6 \\
Tumurkhuu & Veri artefaktları (O1/O2), diller arası aktarım (Tablo~\ref{tab:transfer}); (ii); \S2, \S6 \\
\bottomrule
\end{tabular}
\caption{Görev dağılımı (taslağı herkes okur).}
\label{tab:plan}
\end{table}
